% concatenated from main_3.tex
\documentclass[12pt]{article} 
\usepackage[T2A]{fontenc} 
\usepackage[utf8]{inputenc} 
\usepackage{comment, amsthm, amssymb, amsmath, mathrsfs, bbm, xcolor} 
\usepackage[draft]{hyperref}
\usepackage[english]{babel}
\textheight=22cm 
\textwidth=16cm 
\hoffset=-12mm 
\voffset=-20mm

\newtheorem{conjecture}{Conjecture}
\newtheorem*{id}{Identity}
\newtheorem{theorem}{Theorem} 
\newtheorem{prop}{Proposition}[section]
\newtheorem{lemma}{Lemma}[section]
\newtheorem{cor}{Corollary}[section]
%\theoremstyle{remark}
%\newtheorem{remark}{Remark}

\newtheorem{example}{Example}

\newtheorem{fact}{Fact}

\newtheorem{problem}{Problem}

\usepackage{hyperref}

\newtheorem{solution}{Solution}

\newtheorem{corollary}{Corollary}

\newtheorem{proposition}{Proposition}

\newtheorem*{sketch}{Sketch of the proof}

\newtheorem{rem}{Remark}[section]
%\renewcommand{\therem}{\arabic{rem}}
\theoremstyle{definition}
\theoremstyle{plain}
\newtheorem{definition}{Definition}[section]
\renewcommand{\leq}{\leqslant}
\renewcommand{\geq}{\geqslant}
\renewcommand{\le}{\leqslant}
\renewcommand{\ge}{\geqslant}
\DeclareMathOperator{\End}{End}
\DeclareMathOperator{\Tr}{Tr}
\def\T{\mathcal{T}}

\begin{document}
\newcommand{\eps}{\varepsilon}
\renewcommand{\phi}{\varphi}


\newcommand{\suchthat}{\, : \,}
\newcommand{\commentDK}[1]{{\color{blue}{#1 -- DK}}}
\newcommand{\dima}[1]{\commentDK{#1}}  
\newcommand{\fedya}[1]{{\color{orange}{#1 -- FP}}}
\newcommand{\red}[1]{{\color{red}{#1}}}
\newcommand{\green}[1]{{\color{green}{#1}}}



\author{Ilya I. Bogdanov, Fedor Petrov, Anton Sadovnichiy, Fedor Ushakov}
\title{Biregular bipartite labeled multigraphs and perfect matchings in bipartite tensor products}
\maketitle
\begin{abstract}

In 2019, P. Higgins formulated~\cite{H} a question about bipartite
graphs (see Conjecture \ref{HigginsC} below); this question arises in the study of regular
finite semigroups. F. V. Petrov formulated \cite{P}
another combinatorial conjecture
(Conjecture \ref{BallsC}); Conjecture \ref{BallsC} implies Conjecture \ref{HigginsC} and seems simple itself. However, both conjectures remain
unproven in the general case. In the present paper, some special cases are proved,
 Conjecture \ref{4PartsC} is formulated in the same spirit, and some of its special cases are proved. In addition,
Conjecture \ref{HigginsC} is reduced to a matrix inequality
(Conjecture \ref{MatrixInequalityC}); this inequality is, in turn, also proved in a special case.

\end{abstract}

\section{Introduction}

\begin{definition}
    A bipartite labeled multigraph is a tuple $(L,R,E,u,v)$, where $A$, $B$, and $E$ are finite sets (left part, right part, and the set of edges, respectively); $u$ and $v$ are functions; $u: E \rightarrow L, v : E \rightarrow R$.
\end{definition}

\textbf{Remark.} $L$ and $R$ may intersect.

\begin{definition}
A \emph{perfect matching} in a (multi)graph is an involution $\theta$ on its vertex set $V$ such that $\forall v \in V$ there exists an edge between $v$ and $\theta(v)$. 
\end{definition}





For the $n\times k$ matrix $M$, by $M\otimes M^T$ we denote the symmetric $nk\times nk$ matrix with rows and columns indexed by pairs $(x,y)\in [n] \times [k]$
with entries $(M\otimes M^T)((x_1,y_1),(x_2,y_2)):=M(x_1,y_2)M(x_2,y_1)$ (this is not exactly the tensor
product of matrices, but we still use this notation). The question by Higgins~\cite{H} can be rephrased as follows.

\begin{conjecture}
  \label{HigginsC}
  Let $M$ be a binary $n\times k$ matrix. Assume that the bipartite graph with both parts of size $nk$ with (bipartite) adjacency matrix $M \otimes M^T$ has a perfect matching. Then an ordinary graph with adjacency matrix $M \otimes M^T$ also has a perfect matching.
\end{conjecture}

Conjecture \ref{HigginsC} is equivalent to the following (see Theorem \ref{WeakBallsAndHigginsAreEquiv}):

\begin{conjecture}\label{WeakBallsC}
  Let $M$ be a binary $n\times k$ matrix with non-negative integer entries with row sums $k$ and column sums $n$. Then an (non-bipartite) $nk$-regular multigraph with $nk$ vertices with adjacency matrix $M \otimes M^T$ has a perfect matching.
\end{conjecture}

\begin{definition} \label {SquareBracket}
    For two functions $u$ and $v$ with the same domain $E$, $u : E \rightarrow L, v : E \rightarrow R$, by $[u,v]$, we denote the weighted complete bipartite graph with parts $L$ and $R$, and the weight of the edge $(x,y) \in L \times R$ equal to $| \{ e \in E, u(e) = x, v(e) = y \} |$; we also denote by $[u,v]$ the $L \times R$ matrix with the weights of the edges of this graph. Actually, we identify both objects defined as $[u,v]$.
\end{definition}

Denote by $W_{L,R}$ the weighted complete bipartite graph with parts $L$ and $R$, and with all weights equal to $1$.

\begin{definition} \label {ActionA}
Denote the set of all biregular bipartite labeled multigraphs with $L = [n], R = [k], E=[n] \times [k]$ ($u$ and $v$ can be arbitrary) by $A_{n,k}$;
Denote the group of permutations of $[n] \times [k]$ by $S_{n,k}$.
Consider the following two-sided action of $S_{n,k}$ on $A_{n,k}$: $\sigma (u,v) = (u \circ \sigma^{-1}, v)$, $(u, v) \sigma = (u, v \circ \sigma)$. 
\end{definition}

\begin{definition} \label {ActionB}

For some $L$ and $R$, let $S_L$ be the group of permutations of $L$, and $S_R$ be the group of permutations of $R$. $S_L$ and $S_R$ act on the $L \times R$ -matrices: for a matrix $M$ and $\alpha \in S_L, \beta \in S_R$, $(\alpha M)_{i,j} = M_{\alpha^{-1}(i),j}$; $(\beta M)_{i,j} = M_{i, \beta^{-1}(j)}$. 

\end{definition}

Denote the set of all matrices of type $[n] \times [k]$ with non-negative integer entries, with row sums equal to $k$ and column sums equal to $n$, by $B_{n,k}$;
$[.,.]$ gives a map $A_{n,k} \rightarrow B_{n,k}$. For $G = (u,v) \in A_{n,k}$ we also denote $[u,v]$ by $[G]$

Denote the set of involutions in $S_{n,k}$ by $I = \{ \iota \in S_{n,k}, \iota^2 = id \}$.

The second author has formulated Conjecture \ref{BallsC} in \cite{P}. 

\begin{conjecture}\label{BallsC}
  $\forall (u,v) \in A_{n,k} \exists \iota \in I, [u \circ \iota, v] = W_{[n],[k]}$
\end{conjecture}

In \cite{P}, Conjecture \ref{BallsC} was formulated in terms of the existence of a certain swapping process. Namely, if the parts $L$ and $R$ correspond to
$n$ girls and $k$ colors, we may interpret
every edge $e \in E$ as a ball of the color $v(e)$ held by the girl $u(e)$. Thus, every girl has $k$ balls, and there are $n$ balls of each color. 
The 2-cycle $(e,f), e \neq f, \iota(e)=f$ in $\iota$  corresponds to the swap of the $u(e)$'s ball with color $v(e)$ and the $u(f)$'s ball with color $v(f)$. A fixed point corresponds to a ball that does not participate in a swap.
Thus, Conjecture \ref{BallsC} states that the girls can organize the swapping process so that, ultimately, each girl has exactly one ball of each color and every ball participates in at most one swap.

\medskip
\noindent
\textbf{Remark.}
Conjecture $\ref{BallsC}$ implies Conjecture \ref{WeakBallsC}. Indeed, for $M$ from Conjecture \ref{WeakBallsC}, consider any $(u,v) \in A_{n,k}$ such that $[u,v]'=M$ and $\iota \in I$, $[u \circ \iota, v]=W_{[n],[k]}$. Then $\sigma = (u \circ \iota) \times v$ is a permutation of $[n]\times[k]$, and $\theta = \sigma \circ\iota \circ \sigma^{-1}$ is a perfect matching in $M \otimes M^T$.



\medskip 

\begin{conjecture}\label{4PartsC}
Let $L_1,R_1,L_2,R_2$ be sets of vertices of sizes $n_1,k_1,n_2,k_2$, respectively with $n_1 k_2 = n_2 k_1$, let $(L_1,R_1,E_1,u_1,v_1)$ and $(L_2,R_2,E_2,u_2,v_2)$, where $|E_1|=|E_2|=n_1 k_2$ be biregular bipartite labeled multigraphs;
then there exists a bijection $\psi: E_1 \rightarrow E_2$ such that $[u_2 \circ \psi, v_1] = W_{L_2, R_1}$ and  $[u_1 \circ \psi^{-1}, v_2] = W_{L_1, R_2}$

\end{conjecture}

We have proved its corollary; see Theorem~\ref{Weak4PartsT}.

We also formulate one more conjecture that implies Conjecture \ref{WeakBallsC} (see Theorem~\ref{InequalityImpliesTutte}).

\begin{conjecture}\label{MatrixInequalityC}
Assume that $M$ is a matrix with non-negative real elements, $t$ is an odd positive integer, and $(x_1, y_1), \dots, (x_t, y_t)$ are distinct entries in $M$ such that $M_{x_i y_i} = 0$ for all $i=1,2,\dots,t$. Then 
$$
  \sum_{i=1}^{t} \sum_{j=1}^t M_{x_i y_j} M_{x_j y_i} \leq (t-1)\|M\|_{\infty}\|M\|_{1}.
$$
\end{conjecture} 

(Recall that $\|M\|_\infty$ is the maximum of $\ell^1$-norms of the rows
of $M$, and $\|M\|_1=\|M^T\|_\infty$ is the maximum of $\ell^1$-norms
of the columns of $M$).

\section{Main results}

In the following theorem, the tensor product has the usual meaning. 

\begin{theorem}\label{Weak4PartsT}
Let $n_1, k_1, n_2, k_2$ be positive integers with $n_1 k_2 = n_2 k_1$. Let $M_1$ be an $n_1 \times k_1$ matrix with all the row sums equal to $k_2$, and all the column sums equal to $n_2$; let $M_2$ be an $n_2 \times k_2$ matrix with all the row sums equal to $k_1$, and all the column sums equal to $n_1$. Consider a square binary matrix $M$ with the same nonzero entries as in $M_1 \otimes M_2^T$. Then the bipartite graph with the (bipartite) adjacency matrix $M$ has a perfect matching.
\end{theorem}

\begin{theorem}\label{BallsT}
Conjecture \ref{BallsC} holds for the case $k=mn+\eps, \eps \in \{ -2, -1, 0, 1, 2\}, m \in \mathbb{N}$%, where $n$ and $k$ are the sizes of the parts of the bipartite multigraph
\end{theorem}

Its analog for Conjecture \ref{4PartsC} is

\begin{theorem}\label{4PartsT}
Conjecture~\ref{4PartsC} holds for the case $n_1=n_2=n, k_1=k_2=k, k=mn+\eps, \eps \in \{ -2, -1, 0, 1, 2\}, m \in \mathbb{N}$ and also whenever $k_1$ is divisible by $n_1$.
\end{theorem}

\begin {theorem}\label{InequalityImpliesTutte}
Conjecture $\ref{MatrixInequalityC}$ implies Conjecture $\ref{WeakBallsC}$.
\end {theorem}

\begin{theorem}\label{SquareMatrixT}
Conjecture \ref{MatrixInequalityC} holds for the case when all $x_i$ are distinct and all $y_i$ are distinct.
\end{theorem}

\section{Proof of the main results}

\subsection{Proof of Theorem~\ref{Weak4PartsT}}

In this section, all bipartite graphs are assumed to have nonempty left and right parts.

We can naturally transfer the notion of the tensor product to bipartite (multi)graphs from their bipartite adjacency matrices.

\begin{definition}
\label{h(G)}
Consider the bipartite graph $G$ with parts $(L,R)$. For $X \subseteq L$, by $N(X)$ we denote the union of the neighborhoods of the vertices of $X$. Define $h(G)$ as
$$
  h(G)=\min_{\varnothing \neq X \subseteq L} \frac{|N(X)|}{|X|}.
$$
\end {definition}

\begin {lemma} \label {TensorWithKL}
Let $p,q$ be positive integers, and $G$ be a bipartite graph. Then the following conditions are equivalent:

(1) $h(G) \geq \frac{p}{q}$;

(2) In the graph $G \otimes K_{p,q}$, there exists a matching that covers its left part.
\end {lemma}

\begin {proof}
$(1) \Rightarrow (2)$. Assuming (1), we will prove that the condition of Hall's lemma is satisfied for $G \otimes K_{p,q}$. Denote the left part of $G$ by $L$.

Consider $\varnothing \neq X \subseteq L \times [p]$. Let $X'$ be the image of $X$ in the projection of $L \times [p]$ onto $L$. Then $|X'| \geq \frac{|X|}{p}$. By (1) and Definition~\ref{h(G)}, we have
$$
  \frac{|N(X')|}{|X'|} \geq h(G) \geq \frac{p}{q}.
$$  

Thus $|N(X)| = |N(X') \times [q]| = q|N(X')| \geq p|X'| \geq |X|$, as desired.

$(2)\Rightarrow(1)$. Consider any $\varnothing \neq X \subseteq L$ and put $X' = X \times [p]$. Since there exists a matching in $G \otimes K_{p,q}$ covering its left part, we have $q|N(X)| = |N(X')| \geq |X'| = p|X|$. So, $\frac{|N(X)|}{|X|} \geq \frac{p}{q}$. Since this holds for any nonempty $X \subseteq L$, it follows that $h(G) \geq \frac{p}{q}$.
\end {proof}

\begin {lemma} \label {hIsMultiplicativeL}
For any bipartite graphs $G$ and $H$, $h(G \otimes H) = h(G)h(H)$.
\end {lemma}


\begin {proof} 
We start by proving that $h(G \otimes H) \leq h(G)h(H)$. Let $X$ and $Y$ be subsets of the left parts in~$G$ and~$H$ such that $h(G)=\frac{|N(X)|}{|X|}$ and $h(H)=\frac{|N(Y)|}{|Y|}$. Using $N(X \times Y) = N(X) \times N(Y)$, we see that
$$
  h(G \otimes H) \leq \frac{|N(X \times Y)|}{|X \times Y|} = \frac{|N(X)|}{|X|}\cdot \frac{|N(Y)|}{|Y|} = h(G)h(H).
$$

It remains to show that $h(G \otimes H) \geq h(G)h(H)$. Let $h(G) = \frac{a}{b}$ and $h(H) = \frac{c}{d}$. By Lemma \ref{TensorWithKL}, $G \otimes K_{a,b}$ and $H \otimes K_{c,d}$ have matchings covering their left parts. Then $(G \otimes K_{a,b}) \otimes (H \otimes K_{c,d}) \cong (G \otimes H) \otimes (K_{a,b} \otimes K_{c,d}) \cong (G \otimes H) \otimes K_{ac, bd}$ also has a matching that covers the left part. Using Lemma~\ref{TensorWithKL}, we see that $h(G \otimes H) \geq \frac{ac}{bd} = h(G)h(H)$.
\end {proof}

\begin{proof}[Proof of Theorem \ref{Weak4PartsT}]
Let $G$ and $H$ be the bipartite graphs with bipartite adjacency matrices $M_1$ and~$M_2^T$, respectively. The conditions on the row and column sums yield $h(G) = \frac{k_2}{n_2}$ and $h(H) = \frac{n_1}{k_1} = \frac{n_2}{k_2}$. Thus, by lemma $\ref{hIsMultiplicativeL}$, we have $h(G \otimes H) = 1$. Therefore, there is a perfect matching in $G \otimes H$.
\end {proof}

Consider $u_0:[n]\times[k] \rightarrow [n], u_0((i,j))=i$ and $v_0 : [n] \times[k] \rightarrow [k], v_0((i,j))=j$; 

\begin {theorem} \label{WeakBallsAndHigginsAreEquiv}
  Conjectures \ref{HigginsC} and \ref{WeakBallsC} are equivalent.
\end {theorem}

\begin {proof}

$(1) \Rightarrow (2)$ Assuming Conjecture \ref{HigginsC}, it suffices to check that the bipartite graph with the adjacency matrix $M \otimes M^T$ has a perfect matching. By lemma \ref {hIsMultiplicativeL}, $h(M \otimes M^T) = h(M)h(M^T) = \frac{k}{n} \cdot \frac{n}{k} = 1$.

$(2) \Rightarrow (1)$
 Let $M$ be an $n \times k$ binary matrix such that the bipartite graph with the bipartite adjacency matrix $M \otimes M^T$ has a perfect matching $\theta$ (recall that $\theta$ is an involution on $[n] \times [k]$). Consider $M' = [u_0 \circ \theta, v_0]'$. All non-zero entries in matrix $M'$ are non-zero in $M$. Conjecture \ref{WeakBallsC} implies that the ordinary multigraph with adjacency matrix $M' \otimes M'^T$ has a perfect matching; so does $M \otimes M^T$.

\end {proof}

\subsection{Results related to Conjecture \ref{BallsC}}

We start with the proof of Theorem \ref{BallsT} for the case $n=k$ because it's clearer. We explain it in the girls-and-balls language. Recall that K\"onig's theorem states that every regular bipartite multigraph (not necessarily a graph) with positive degrees has a perfect matching. (see \cite{K})

Consider the $n\times n$ board with rows corresponding to the girls.
By K\"onig's theorem, there exist balls
$Ball_1,\ldots,Ball_n$ held by girls 1,2,\ldots, $n$, respectively, of mutually different colors. Put these $n$ balls in the corresponding cells of
the first column of our board. Removing these $n$ balls, we still have a regular
multigraph (with degree $n-1$ for every girl and every color). Do the same and fill
the second column of our board. Proceeding in this way, we fill the board with the balls. Note that the transposition of the board corresponds to a series of disjoint swaps (and exactly $n$ balls, one per girl, remain unswapped). In the end, the $i$-th girl has the balls that were placed in the $i$-th column of our board; thus, all are of different colors. 
\qed

\begin{lemma} \label{Coloring}
For any bipartite multigraph $G$ and any positive integer $m$, we can color the edges of this multigraph in $m$ colors such that for every vertex and for every two colors, the degrees of this vertex with respect to these two colors differ by at most $1$.
\end{lemma}

\begin{proof} We start by proving the lemma statement for the case $m=2$. Call these two colors Blue and Red. Add one new vertex to each part of $G$ and connect them to some old vertices and possibly to each other to make all degrees even. Now the edges of the graph can be divided into even cycles. Let us traverse them, coloring the edges alternately blue and red; we obtain the coloring of all the edges of the new graph such that each vertex has equal red and blue degrees.
By removing new edges and new vertices, we obtain a coloring that satisfies the condition.

For an arbitrary coloring of the edges of~$G$, define its \emph{cost} as the sum of the squares of the colored degrees over all the vertices over all colors. The sum of squares of the two colored degrees of a certain vertex is minimized if and only if those degrees differ by at most~$1$. Therefore, the minimum value of the cost is achieved exactly when the coloring satisfies the requirements of the Lemma.

Now we proceed with the general case of an arbitrary $m$. Choose a coloring of minimum cost. Suppose that for some two colors, say Red and Blue, there is a vertex whose red and blue degrees differ by more than 1. Consider the subgraph containing only red and blue edges; as we have shown, its edges can be recolored in red and blue so that the blue and red degrees of each vertex differ by at most $1$. As mentioned above, this decreases the cost of the whole coloring, which leads to a contradiction. Therefore, the chosen coloring satisfies the requirements.
\end{proof}

\begin{lemma} \label{windFromColoring}
If $(L,R,E,u,v)$ is a biregular bipartite labeled multigraph, $|L|=n,|R|=k,|E|=nk,k=mn+\eps, |\eps| < n$, $F$ is a set of size $n$, there exists a coloring $w : E \rightarrow F$ such that $[w,v]=W_{F,R}$ and all weights in $[u,w]$ are equal to $m$ or $m+{\rm sgn}\,(\eps)$.
\end{lemma}

\begin {proof}
We can use $F$ as the set of colors for Lemma \ref{Coloring} to obtain a coloring function $w: E \rightarrow F$; The weights of the edges in $[v,w]$ correspond to the colored degrees in this coloring for the vertices of $R$. All vertices in~$R$ have all colored degrees equal to $1$, so $[v,w]=W_{F,R}$.
\end {proof}

\textbf{Remark.}
When applying Lemma \ref{windFromColoring}, we call \emph{exceptional} the edges of $[u,w]$ with weight $m+{\rm sgn}\,(\eps)$. The bipartite graph of exceptional edges is $|\eps|$-regular.

%\begin {definition}
%We call a bipartite labeled multigraph $H=(L,R,E,u,w)$ $\emph{symmetrizable}$ if there exists a bijection $\phi : L \rightarrow R$ and an involution $\iota$ on $E$ such that $u \circ \iota = \phi^{-1} \circ w$.

%We call a weighted complete bipartite graph with parts $L$ and $R$ $\emph{symmetrizable}$ if there exists a bijection $\phi : L \rightarrow R$ such that $\forall x \in L, y \in R$, the weight of the edge $(x,y)$ equals to the weight of the edge $(\phi^{-1}(y), \phi(x))$.

%We call a bipartite graph symmetrizable if the corresponding weighted bipartite graph with weights 0 and 1 is symmetrizable.

%\end {definition}

%\begin {prop} \label {xiSymm}
%A bipartite labeled multigraph $H$ is symmetrizable iff $\chi(H)$ is.
%\end {prop}

%\begin {prop} \label {regularSymmetrizable}
%Any 0-, 1-, or 2- regular bipartite graph is symmetrizable.
%\end {prop}

%\begin{lemma} \label{WeNeedSymmetrizablewind}

%Conjecture \ref{BallsC} holds for an $n,k$-wreath $G$ iff $\exists$ symmetrizable wind $H=(L,L,[n]\times[k],u,w)$ that matches $G$.

%\end{lemma}

%\begin {proof}

%$\Rightarrow$ Consider $\iota$, $\chi(\iota \circ u, v) = W_{[n],[k]}$. The wind $H = (L, L, [n] \times [k], u, u \circ \iota)$ matches $G$ and is simmetrizable with $\phi = id_L$ and $\iota$.

%$\Leftarrow$ $\chi(\iota G) = \chi(u \circ \iota, v) = \chi(\phi^{-1} \circ w, v) = W_{[n],[k]}$. 

%\end {proof}

\begin {proof}[Proof of Theorem \ref{BallsT}.]
Take any coloring $w : [n] \times [k] \rightarrow [n]$ from Lemma \ref{windFromColoring}. There exists an involution $\iota$ on $[n] \times [k]$ and $\alpha \in S_n$ such that $u \circ \iota = \alpha \circ w$. Thus, $[u \circ \iota, v] = [\alpha \circ w, v] = \alpha [w,v] = \alpha W_{[n],[k]} = W_{[n],[k]}$.
\end {proof}

For $n=6$, the only remaining case is considered in the following theorem.

In a bipartite multigraph with left part of size $6$ and right part of size $6m+3$, where all vertices in the left part have multidegree $2m+1$ and all vertices in the right part have multidegree $2$, we say that two distinct vertices $u$ and $v$ from the left part are \emph{complementary} if $|N(\{u,v\})| = 2m+1$. Note that each vertex in the left part has at most one complementary vertex.

\begin{lemma}\label{find3}
For a bipartite multigraph as described above, let $u$, $v$ and $f$ be three vertices from its left part such that $u$ and $v$ are not complementary. Then there exists a vertex $w$ in the left part, $w \notin \{u,v,f\}$, such that $|N(\{u,v,w\})| \ge 3m+3$ and no two vertices among $\{u,v,w\}$ are complementary.
\end{lemma}

\begin{proof}
Let $l = |N(\{u,v\})|$. We consider several cases.

\emph{Case 1:} There exist complementary vertices among the four vertices of the left part different from $u$ and $v$ (i.e., among the set consisting of $f$ and the three other vertices). Choose such a complementary pair and let $w$ be the vertex of that pair that is not $f$. Then $w$ is not complementary to $u$ or $v$ because each vertex has at most one complementary vertex. Moreover,
\[
|N(\{u,v,w\})| = |N(\{u,v\})| + |N(w)| \ge (2m+2) + (m+1) = 3m+3.
\]

\emph{Case 2:} Both $u$ and $v$ have complementary vertices (say $u'$ and $v'$, respectively). Then we can choose $w$ to be any vertex distinct from $u,v,f$ that is not complementary to $u$ or $v$ (such a vertex exists because there are six vertices in total and at most two complements). Then
\[
|N(\{u,v,w\})| \ge |N(u)| + |N(v)| + |N(w)| \ge (m+1) + (m+1) + (m+1) = 3m+3.
\]

\emph{Case 3:} Only one of $u$ and $v$, say $u$, has a complementary vertex. Then we can take two different vertices $w_1$ and $w_2$ from the left part, different from $u,v,f$, that are not complementary to $u$. Note that $l = |N(\{u,v\})| \ge (m+1) + |N(v)|$. The total number of edges from $w_1$ and $w_2$ to $N(\{u,v\})$ is at most $2|N(v)| - (2m+1) < 2(2|N(v)| - 2m - 1) \le 2(2l - 4m - 3)$. Hence for at least one of $w_1$ and $w_2$, say $w_1$, the number of edges to $N(\{u,v\})$ is less than $2l - 4m - 3$, so the number of edges to the remaining vertices of the right part is greater than $6m+4 - 2l$. Consequently, $|N(\{u,v,w_1\})| \ge 3m+3$.

\emph{Case 4:} Neither $u$ nor $v$ has a complementary vertex. Observe that there are at most two vertices $s$ different from $u$ and $v$ such that $|N(\{u,v,s\})| < 3m+3$. Therefore we can choose $w$ among the remaining vertices (excluding $f$) satisfying the required condition.
\end{proof}

\begin{theorem}
Conjecture \ref{BallsC} holds for the case $n=6$, $k = 6m+3$.
\end{theorem}

\begin{proof}
First, we color the edges into three colors by Lemma \ref{Coloring}. For each of the three colors we aim to find three vertices $x,y,z$ such that none of them are complementary in the multigraph of that color and $|N(\{x,y,z\})| \ge 3m+3$. Then we can recolor the edges of that color into two colors so that in the first of these two colors the vertices $x,y,z$ have degree $m+1$, all other left vertices have degree $m$, and every right vertex has degree $1$.

Take any two non‑complementary vertices in the multigraph of the first color, and let $f$ be any other vertex from the left part. Set $u$ and $v$ as these two non‑complementary vertices. By Lemma \ref{find3} we obtain a vertex $w$ (distinct from $u,v,f$) such that $u,v,w$ satisfy the required conditions. This gives the desired triple for the first color.

Thus we have three vertices $1,b,c$ with $|N(\{1,b,c\})| \ge 3m+3$ in the first color. Without loss of generality we may assume that $b$ and $c$ are not complementary in the second color. For the multigraph of the second color, take $1,b,c$ as $f,u,v$ respectively and apply Lemma \ref{find3} to find another vertex; call it $4$. Then $b,c,4$ form the required triple for the second color.

Now rename $b$ and $c$ as $2$ and $3$ so that $3$ and $4$ are not complementary in the third color. For the third color, take $f,u,v$ to be $2,3,4$ respectively and apply Lemma \ref{find3} again. The vertex obtained may be $1$ or a new vertex, say $5$.

If it is $5$, we have the triples $(1,2,3)$, $(2,3,4)$ and $(3,4,5)$ for the first, second and third colors respectively.  
If it is $1$, we instead apply Lemma \ref{find3} to the first color with $f,u,v = 1,2,3$ to find either $4$ or $5$ and take this new triple for the first color. Consequently we obtain either the triples $(2,3,4)$, $(2,3,4)$, $(3,4,1)$ or the triples $(2,3,5)$, $(2,3,4)$, $(3,4,1)$ for the three colors. In any case, by the same argument as in Theorem \ref{BallsT}, we can find an involution.
\end{proof}

\subsection{Results related to Conjecture ~\ref{4PartsC}}

First, we present the observation that motivated Conjecture $\ref{4PartsC}$.

\begin {theorem}
If Conjecture $\ref{4PartsC}$ is true and Conjecture $\ref{BallsC}$ is true for connected multigraphs, then Conjecture $\ref{BallsC}$ is true.
\end {theorem}

\begin {proof}
Suppose there exists the smallest disconnected $G \in A_{n,k}$ such that Conjecture \ref{BallsC} fails for $G$. Suppose $G = G_1 \sqcup G_2, G_1 \in A_{n_1,k_1}, G_2 \in A_{n_2,k_2}$, $g=gcd(n,k)$. By Lemma \ref {Coloring}, we can color the edges of $G$ in $g$ colors so that all colored degrees are the same for every vertex. Take $t_i = g n_i/n$. Divide the edges of $G_i$ into two biregular bipartite labeled multigraphs with the same parts: edges of the first $t_i$ colors go to $T_i$, and all the other edges go to $S_i$. Conjecture $\ref{BallsC}$ must hold for $T_i$, and Conjecture $\ref{4PartsC}$ holds for the pair of $S_1$ and $S_2$; thus, Conjecture $\ref{BallsC}$ holds for $G$, which is a contradiction.
\end {proof}

\begin{lemma} \label{makeHamiltonian}
Let $k= mn+\eps, \eps \in \{ -2, 2\}, m \in \mathbb{N}, 
n>2$. For a biregular bipartite labeled multigraph $(L,R,E,u,v)$ with $|L_1|=n,|R|=k,|E|=nk$, and a set $F, |F|=n$, there exists a coloring $w : E \rightarrow F$ as in Lemma~\ref{windFromColoring} such that the graph of exceptional edges in $[u,w]$ is connected (and thus it is just a cycle of length $2n$).
\end{lemma}


\begin{proof}
Choose $w$ as in Lemma \ref{windFromColoring}, minimizing the number of connected components in the graph of exceptional edges in $[u,w]$. Denote this graph by $C$. Suppose $C$ is disconnected. Let Red and Blue $\in F$ be two colors from different components. Leave only blue and red edges in $G$ and orient the red edges from $L$ to $R$ and the blue edges from $R$ to $L$ to obtain the digraph $D$. We call a vertex of $L$ positive if its outdegree in $D$ is greater than its indegree, and negative if it is less. Denote one of the two neighbors of Red in $C$ by $s$. Since Red and Blue are in different components of $C$, $s$ is not joined with Blue in $C$; i.e., the blue degree of $s$ is $m$. Therefore, $s$ has the same sign as $\eps$.

Assume that $\eps>0$.
Denote the set of vertices reachable from $s$ in $D$ by $S$.
Note that there are no edges from $S$ outwards (but there can be edges inwards). 
Then, by the indegree/outdegree balance for the induced subgraph of $D$ on $S$, it is impossible for all vertices in $S$ to
be nonnegative (since at least one of them, namely, $s$, is positive). In other words, there exists a path from $s$ to a negative vertex $t$. Flip the red and blue color of the edges in this path. 
In the graph $C$, this corresponds to the replacement of edges $s$-Red and $t$-Blue with $s$-Blue and $t$-Red. For a regular graph of degree 2 (a disjoint union of cycles), this operation decreases the number of connected components. A contradiction. 
Thus, $C$ is connected. The case $\eps<0$ is analogous. 
\end{proof}

%\begin {definition} \label{Isomorphicwinds}

%We say that biregular bipartite labeled multigraphs $H_1 = (L_1, L_2, E_1, u_1, w_1)$ and $H_2 = (L_2, L_1, E_2, u_2, w_2)$ are \emph{isomorphic} if there exists a permutation $\sigma_1$ of $L_1$ and $\sigma_2$ of $L_2$, and a bijection $\psi : E_1 \rightarrow E_2$ such that $u_2 \circ \psi = \sigma_2 \circ w_1$ and $u_1 \circ \psi^{-1} = \sigma_1 \circ w_2$.

%\end {definition}

%\begin {prop} \label {xiIsomorphic}

%$H_1$ and $H_2$ are isomorphic iff $[u_1, w_1]$ and $[u_2, w_2]$ are isomorphic in common sense with fixed parts.

%\end {prop}

%\begin{lemma} \label{WeNeedIsomorphicwinds}

%Conjecture \ref{4PartsC} holds for the winds $G_1$ and $G_2$ iff there exist isomorphic winds $H_1$ and $H_2$ such that $H_i$ matches $G_i$.

%\end{lemma}


%\begin{proof}

%$\Rightarrow$ Consider $H_1 = (L_1, L_2, E_1, u_1, u_2 \circ \psi), H_2 = (L_2, L_1, E_2, u_2, u_1 \circ \psi^{-1})$; they match with $G_1$ and $G_2$ respectively by the definition of $\psi$ from Conjecture \ref{4PartsC}. $H_1$ and $H_2$ are isomorphic with $\sigma_i = id_{L_i}$.

%$\Leftarrow$ $[u_2 \circ \psi, v_1] = [\sigma_2 \circ w_1, v_1] = [w_1, v_1] = W_{L_2,R_1}$.  Analogously, $[u_1 \circ \psi^{-1}, v_2] = W_{L_1, R_2}$. 
%\end{proof}

\begin{proof}[Proof of Theorem \ref{4PartsT}]
 If $k_1$ is divisible by $n_1$, say $k_1 = m n_1$, then $k_2 = m n_2$ and by Lemma \ref{windFromColoring}, there exist $w_1 : E_1 \rightarrow L_2$ and $w_2 : E_2 \rightarrow L_1$ such that $[w_1,v_1]=W_{L_2,R_1}$ and $[w_2,v_2]=W_{L_1,R_2}$ and $[u_i,w_i]$ have no exceptional edges.
 
Else if $k = mn+ \eps, \eps \in \{-1, 1\}$, then by Lemma \ref{windFromColoring} there exist analogous $w_1$ and $w_2$ such that the graphs of exceptional edges in $[u_i, w_i]$ are $1$-regular.

Else $n, k$ satisfy the condition of Lemma \ref{makeHamiltonian}, so there exist $w_1$ and $w_2$ such that the graphs of exceptional edges in $[u_i,w_i]$ are cycles of length $2n$. 

Anyway, we obtain colorings $w_1,w_2$ such that $[w_1,v_1]=W_{L_2,R_1}$ and $[w_2,v_2]=W_{L_1,R_2}$ and $[u_i,w_i]$ are isomorphic with fixed parts. This isomoprhism implies that there exist permutations $\sigma_i$ of $L_i$ and a bijection $\psi : E_1 \rightarrow E_2$ such that $u_2 \circ \psi = \sigma_2 \circ w_1$ and $u_1 \circ \psi^{-1} = \sigma_1 \circ w_2$. Thus, $[u_2 \circ \psi, v_1] = [\sigma_2 \circ w_1, v_1] = \sigma_2 [w_1, v_1] = \sigma_2 W_{L_2, R_1} = W_{L_2, R_1}$. Analogously, $[u_1 \circ \psi^{-1}, v_2] = W_{L_1, R_2}$.
\end{proof}


\subsection{Proof of Theorem~\ref{InequalityImpliesTutte}}

We need a version of Tutte's theorem on perfect matchings (see \cite{T}) for graphs with loops

\begin{lemma}\label{TutteWithLoops}
Let $G=(V,E)$ be a graph, possibly with loops. Then $G$ has a perfect matching (where loops are allowed) if and only if, 
for every subset $U
\subset V$, the graph $G-U$ has at most
$|U|$  odd connected components without loops.
\end{lemma}


\begin{proof}
If a perfect matching exists, then for every odd connected component of $G-U$ without loops, it must contain an edge from
this component to $U$; thus, the condition is satisfied. 

Conversely, suppose the condition is satisfied. 
Denote the set of vertices with loops by $L$; $l=|L|$.
Add $l$ or $l+1$ new vertices so that the total number of vertices becomes even. Connect all the new vertices with each other and with every vertex from $L$, and remove all the loops. We obtain a graph $G'$ without loops. If $G'$ contains a perfect matching, it naturally corresponds to a perfect matching in $G$. 

Suppose $G$ has no perfect matching. Then $G'$ also has no perfect matching; thus, by Tutte theorem,  
there exists a set of vertices $A$ in $G'$
such that $G'-A$ contains at least $|A|+1$ odd components.
By a parity argument, there are actually more than $|A|+1$ odd components.
If $A$ contains some of the new vertices but not all, then if we remove all new vertices from $A$, the number of odd components in $G'-A$ will decrease by at most $1$, and the size of $A$ will decrease by at least one.
If $A$ contains all of the new vertices, then if we remove all new vertices from $A$, the number of odd components in $G'-A$ will decrease by at most $l$, and the size of $A$ will decrease by at least $l$. 

So we can assume that $A$ doesn't contain any new vertices. Then all vertices from $L$ are in one component in $G'-A$. Therefore, $G-A$ has more than $|A|$ odd components without loops, a contradiction.
\end{proof}




\begin {proof}[Proof of Theorem \ref{InequalityImpliesTutte}]
 Consider the matrix $M$ from the statement of Conjecture $\ref{WeakBallsC}$. It suffices to verify the condition from lemma \ref{TutteWithLoops} for the multigraph $M \otimes M^T$. The vertices of the graph $M \otimes M^T=(V,E)$ correspond to the cells of the matrix $M$, and the vertices without loops correspond to the cells with zeros. The multigraph $M \otimes M^T$ is $nk$-regular, so, it suffices to verify that for every set $V_1\subset V$
of odd size without loops,
at least $nk$ edges go out from $V_1$ (this would imply
that for every, say, $m$ disjoint odd sets in $V$ without loops,
these edges must go out to at least $m$ distinct vertices, that is the condition of Lemma \ref{TutteWithLoops}.) 

Consider the set $V_1$ consisting of $t$ vertices corresponding to the cells $(x_1, y_1), \dots , (x_t, y_t)$, where $t$ is odd and $M_{x_i y_i} = 0$
(in other words, $V_1$ does not contain loops). The edge multiplicity between $(x_i,y_i)$ and $(x_j,y_j)$ is $M_{x_i y_j}M_{x_j y_i}$. Hence the number of edges outward from $V_1$ is $nkt - \sum_{i=1}^{t} \sum_{j=1}^t M_{x_i y_j} M_{x_j y_i}$, so it suffices to prove that $\sum_{i=1}^{t} \sum_{j=1}^t M_{x_i y_j} M_{x_j y_i} \leq (t-1)nk$. It follows from Conjecture $\ref{MatrixInequalityC}$.
\end {proof}

\subsection{Several proofs of Theorem~\ref{SquareMatrixT}}

\begin{lemma}\label{LambdasSum0}
If $t$ is odd, $\lambda_1, \dots ,\lambda_t \in [-\kappa, \kappa], \sum_{i=1}^t \lambda_i = 0$, then $\sum_{i=1}^t \lambda_i^2 \leq (t-1)\kappa^2$.
\end{lemma}

\begin{proof}
Consider $K \subset \mathbb{R}^t, K = \{ (\lambda_1, \dots, \lambda_t), \lambda_1, \dots ,\lambda_t \in [-\kappa, \kappa], \sum_{i=1}^t \lambda_i = 0 \}$. Since $K$ is compact, the function $\sum_{i=1}^t \lambda_i^2$ attains a maximum on $K$. Let $(\lambda_1, \dots, \lambda_t)$ be the maximum point. Then, for any $i \neq j$, $|\lambda_i| = \kappa$, or $|\lambda_j| = \kappa$, we can move 
the numbers $\lambda_i,\lambda_j$ away from each other while preserving the sum and increasing the sum of squares. Thus, all $\lambda_i$, except possibly one, are equal to $\kappa$ or $-\kappa$. Let us choose $i$ such that $\forall j \neq i$, $\lambda_j \in \{-\kappa, \kappa \}$. $\sum_{i=1}^t \lambda_i = 0$ and $|\lambda_i| \leq \kappa$. Therefore, among $\lambda_j, j \neq i$, the number of $\kappa$ and $-\kappa$ is the same. Hence $\lambda_i = 0$, $\sum_{i=1}^t \lambda_i^2 = (t-1)\kappa^2$
\end {proof}

\begin{proof}[Geometric proof of Theorem \ref{SquareMatrixT}]
Consider the $t \times t$ matrix $B_{ij} = (M_{x_iy_j}M_{x_jy_i})^{1/2}$.

Let $\lambda_1, \dots, \lambda_t$ be the eigenvalues of $B$. Since $B$ is symmetric, all $\lambda_i$ are real.

We have $\sum_{i=1}^{t} \sum_{j=1}^t M_{x_i y_j} M_{x_j y_i}= \Tr B^2= \sum_{i=1}^{t} \lambda_i^2$.

Let $B_i$ be the $i$-th row of the matrix $B$. Then $\|B_i\|_1 = \sum_{j=1}^t B_{ij} = \sum_{j=1}^t (M_{x_iy_j}M_{x_jy_i})^{1/2} \leq (\sum_{j=1}^t M_{x_iy_j}\sum_{j=1}^t M_{x_jy_i})^{1/2} \leq (\|M\|_{\infty} \|M^T\|_{\infty})^{1/2}$.

Hence $|\lambda_i| \leq \|B\|_{\infty} \leq (\|M\|_{\infty} \|M^T\|_{\infty})^{1/2}$ for all $i$.

Moreover, since $B_{ii} = 0$, it follows that $\sum_{i=1}^t \lambda_i = 0$, and by Lemma
\ref{LambdasSum0} $\sum_{i=1}^{t} \lambda_i^2 \leq (t-1)\|M\|_{\infty} \|M^T\|_{\infty}$.

\end{proof}

\begin{proof}[Algebraic proof of Theorem \ref{SquareMatrixT}]
Consider the singular skew-symmetric matrix 
$$
S_{ij}=\begin{cases}
			(M_{x_iy_j}M_{x_jy_i})^{1/2}, & \text{ $i < j$ }\\
            -(M_{x_iy_j}M_{x_jy_i})^{1/2}, & \text{otherwise}
		 \end{cases}
$$

$\Tr S^2 = - \sum_{i=1}^{t} \sum_{j=1}^t M_{x_i y_j} M_{x_j y_i}$.

Denote the eigenvalues of $S$ by $\mu_1, \dots, \mu_t$. Analogously to the bound for $\lambda_i$ from the previous proof, $|\mu_i| \leq (\|M\|_{\infty} \|M^T\|_{\infty})^{1/2}$. Moreover, $det(S)=0$, so one of $\mu_i$ is zero, thus $|\Tr S^2| \leq \sum_{i=1}^{t} |\mu_i|^2 \leq (t -1)\|M\|_{\infty} \|M^T\|_{\infty}$.
\end{proof}

In the following two proofs, without loss of generality, $x_i=y_i=i$. Also, we may assume that $\|M\|_\infty=1\leqslant \|M\|_1$.
Then it suffices to prove that $F:=\sum_{i<j} M_{ij}M_{ji}\leqslant (t-1)/2$ whenever
$M_{ij}$ are non-negative real numbers with restrictions $M_{ii}=0$ and
$\sum_j M_{ij}\leqslant 1$ for all $i$.

\begin{proof}[Combinatorial proof of Theorem \ref{SquareMatrixT}]
Fix all $M_{ij}$ except the first row of matrix $M$ (i.e., for all $i>1$ and all $j$). Then $F$ is multi-affine with respect to $M_{12},\ldots,M_{1t}$,
and its maximal value under restrictions $\sum_j M_{1j}\leqslant 1$ is obtained when one of
$M_{1j}$ equals 1 and others are equal to 0. Change the first row accordingly. Now do the same with the second row, etc. Finally, we have exactly one element 1 in every row. At most $(t-1)/2$ of them form
symmetric pairs, thus indeed $\sum_{i<j} M_{ij}M_{ji}\leqslant (t-1)/2$. 
\end{proof}

The next argument is attributed to Maxim Erogov.

\begin{proof}[Probabilistic proof of Theorem \ref{SquareMatrixT}]
Let number $i$ love number $j$ with probability $M_{ij}$
(these events are disjoint by $j$ and independent by $i$, with some non-negative probability $i$ does not love nobody). Then 
$\sum_{i<j} M_{ij}M_{ji}$ is the expected number of pairs of numbers
mutually loving each other. Obviously, there are always at most $(t-1)/2$
such pairs, thus the result. 
\end{proof}

\section{Further discussion}

Conjecture $\ref{BallsC}$ can also be reformulated in group-theoretic language.

We denote by $S_k^n$ and $S_n^k$ the left and right stabilizers of $(u_0,v_0)$, respectively.

The normalizers of $S_k^n$ and $S_n^k$ in $S_{n,k}$ are $S_k \wr S_n$ and $S_n \wr S_k$, respectively.

\begin {theorem}
The following are equivalent: 

(1) Conjecture \ref{BallsC}; 

(2) $IS_n^kS_k^n=S_{n,k}$; 

(3) $I (S_n \wr S_k) (S_k \wr S_n) = S_{n,k}$.

\end {theorem}

\begin {proof}

(1) $\Rightarrow$ (2): If Conjecture \ref{BallsC} is true, then $\forall \sigma \in S_{[n],[k]}$, $\exists \iota \in I, [ \iota \sigma (u_0, v_0)] = W_{[n],[k]}$; so $\exists \xi \in S_n^k, \xi^{-1} \iota \sigma (u_0, v_0) \xi = (u_0, v_0)$. Thus, $\sigma = \iota \xi (\xi^{-1} \iota \sigma) \in I S_n^k S_k^n$.

(2) $\Rightarrow$ (3) is trivial.

(3) $\Rightarrow$ (1): We need to find $\iota \in I$ such that $[\iota \sigma (u_0, v_0) \pi] = W_{[n],[k]}$ for arbitrary $\sigma, \pi \in S_{n,k}$. Since $\pi^{-1} I \pi = I$, it is equivalent to finding $\iota$, $[\iota \pi \sigma (u_0, v_0)] = W_{[n],[k]}$. Take $\iota \in I, \xi \in S_n \wr S_k, \eta \in S_k \wr S_n$, $\pi \sigma = \iota \xi \eta$. Then $[\iota \pi \sigma (u_0, v_0)] = [\eta (u_0, v_0) \xi] = \alpha \beta [\eta' (u_0, v_0) \xi']$ for some $\alpha \in S_n, \beta \in S_k$, $\eta' \in S_k^n$, $\xi' \in S_n^k$.

\end {proof}

\section{Acknowledgements}

We are grateful to the Adygeya problem solving workshops of 2019
and 2024 organized by the Combinatorics group in MIPT, where a significant part of the progress on 
the topic was achieved. Namely, the case $n=k$ was 
proved by the first author at the 2019 workshop, and later at the
same workshop, the joint efforts of him, Rom Pinchasi, and Ran Ziv allowed for the 
establishment of the $\eps \in \{-1,0,1\}$ case of Theorem \ref{BallsT}. Later,
the $|k-n|\leqslant 2$ case of  Theorem \ref{BallsT}
was obtained by a different argument by Narek Oganisyan
from St. Petersburg University. 
Theorem \ref{4PartsT} was obtained at the 2024 Workshop
by the third and fourth authors.
We would like to thank Sergey Onishchenko from St. Petersburg State University for pointing out a version of Tutte's theorem on perfect matchings for graphs with loops.

The work of F. Petrov and F. Ushakov was performed at the Saint Petersburg Leonhard Euler
International Mathematical Institute and was supported by the Ministry of Science and Higher
Education of the Russian Federation (agreement no. 075–15–2025–343).


%%==================================================================
\begin{thebibliography}{9}%% Параметр определяет ширину номера, разумные значения при нумерации числами: 9 или 99


\bibitem{H} P. Higgins. Permutations and involutions in binary arrays, https://mathoverflow.net/q/324202

\bibitem{P} F. Petrov. Resolution of multiple edges, https://mathoverflow.net/q/325477

\bibitem{K} R. Rizzi, "König's edge coloring theorem without augmenting paths", Journal of Graph Theory, 1998

\bibitem{T} W. T. Tutte, "The factorization of linear graphs", Journal of the London Mathematical Society, 1947

\end{thebibliography}
%%==================================================================
\end{document}
%%==================================================================
